%! Quadratic Functions Revisited — Unit 2, Lesson 2.1 — Warm-Up
\documentclass[10pt]{article}
\usepackage{precalculus-article}
\usepackage{precalculus-boxes}
\newcommand{\TallMath}[1]{$\displaystyle #1\rule[-1.4em]{0pt}{3.2em}$}

\begin{document}

\pageheader{Unit 2, Lesson 2.1}{Warm-Up}

\vspace{0.12in}

\begin{enumerate}[itemsep=10pt]

\item From Lesson 1.3: find the slope of the line through $(2, 7)$ and
      $(6, 19)$.

\begin{work}
  m &= \frac{19 - 7}{6 - 2} \\
    &= \frac{12}{4} \\
    &= 3
\end{work}

$m = $ \blank{3cm}

\item From Lesson 1.1: let $f(x) = x^2 - 4x + 3$. Complete the table. (Watch
      the parentheses: $f(3) = (3)^2 - 4(3) + 3$.)

\smallskip
\renewcommand{\arraystretch}{1.5}
\begin{tabular}{|c|c|c|c|c|c|}
\hline
$x$ & 0 & 1 & 2 & 3 & 4 \\
\hline
$f(x)$ & \blank{1cm} & \blank{1cm} & \blank{1cm} & \blank{1cm} & \blank{1cm} \\
\hline
\end{tabular}
\smallskip

\item The table below gives a function $g$ at equally spaced inputs. A
      \textbf{first difference} is one output minus the output before it.
      Fill in the row of first differences --- no explanation needed yet.

\smallskip
\renewcommand{\arraystretch}{1.5}
\begin{tabular}{|c|c|c|c|c|c|}
\hline
$x$ & 0 & 1 & 2 & 3 & 4 \\
\hline
$g(x)$ & 2 & 5 & 10 & 17 & 26 \\
\hline
\end{tabular}
\smallskip

First differences: \quad
\blank{1.2cm} , \blank{1.2cm} , \blank{1.2cm} , \blank{1.2cm}

\end{enumerate}

\end{document}
